\section{Moment coordinates on the whole space}\label{sec:moment}

The variance argument uses a change of variables whose Jacobian is a
positive Hessian.  The following lemma supplies this change of variables
and the uniform upper bound needed for the later cutoff arguments.

\begin{lemma}[Moment coordinates]\label{lem:moment}
Let $n\ge2$, $\eps>0$, and let $V\in C^\infty(\R^n)$ be even with
$D^2V\ge\eps I$.  Set
\[
 Z=\int_{\R^n}e^{-V(x)}\,dx,
 \qquad d\mu(x)=Z^{-1}e^{-V(x)}\,dx.
\]
There is an even $\varphi\in C^\infty(\R^n)$ such that
$d\nu(z)=e^{-\varphi(z)}\,dz$ is a probability measure,
$\nabla\varphi$ is a diffeomorphism of $\R^n$ onto $\R^n$, and
\begin{equation}\label{eq:moment}
 (\nabla\varphi)_\#\nu=\mu,
 \qquad 0<D^2\varphi\le\frac4\eps I,
 \qquad
 \log\det D^2\varphi
   =V(\nabla\varphi)+\log Z-\varphi.
\end{equation}
Moreover, $\varphi(z)\ge c|z|-b$ for some $c>0$ and $b<\infty$.
The lower Hessian bound in \eqref{eq:moment} means positive definiteness
at every point; it need not be uniform on $\R^n$.
\end{lemma}

\begin{proof}
Evenness and uniform convexity give
$V(x)\ge V(0)+\eps|x|^2/2$, so $0<Z<\infty$.
For $R\ge2$, write $B_R=\{x:|x|<R\}$ and set
\[
 Z_R=\int_{B_R}e^{-V(x)}\,dx,
 \qquad d\mu_R=Z_R^{-1}\mathbf1_{B_R}e^{-V(x)}\,dx.
\]
We first use two compact-support results.  The moment theorem of
Berman--Berndtsson \cite[Theorem~1.1]{BermanBerndtsson2013}, applied to
the body $\overline B_R$ and the positive smooth density
$e^{-V}/Z_R$, gives a smooth convex potential $\varphi_R$, unique up to
translation, for which $\nabla\varphi_R:\R^n\to B_R$ is a
diffeomorphism and
\begin{equation}\label{eq:moment-truncated}
 \det D^2\varphi_R
   =Z_R\exp\{V(\nabla\varphi_R)-\varphi_R\}.
\end{equation}
The centering hypothesis holds by evenness.  Integrating the equation
through this diffeomorphism shows that
$d\nu_R=e^{-\varphi_R}\,dz$ is a probability and
$(\nabla\varphi_R)_\#\nu_R=\mu_R$.
Klartag's estimate \cite[Proposition~3.1 and Remark~3.5]{KlartagArxiv2013}
gives
\[
 0<D^2\varphi_R\le CI,\qquad C=4/\eps.
\]
Indeed, the ball has smooth positively curved boundary, and
$\rho_R=V+\log Z_R$ is convex and smooth, with itself and all its
derivatives bounded on $B_R$, and $D^2\rho_R\ge\eps I$.
The constant $C$ is independent of $R$.

Translate the unique preimage of $0$ to $0$.  It is the unique minimum
of $\varphi_R$.  Reflection gives a second potential for the same
target with its minimum at $0$; uniqueness up to translation therefore
makes $\varphi_R$ even.  For each fixed $R$, take the preimages of
the slopes $\pm(R/2)e_i$.  Their supporting planes give
\[
 \varphi_R(z)\ge\frac{R}{2\sqrt n}|z|-b_R
\]
with a finite, possibly $R$-dependent constant $b_R$.
Thus $\nu_R$ has an exponential tail.  In particular, since
$|\nabla\varphi_R|\le R$, integration of
$\operatorname{div}(z e^{-\varphi_R})$ with expanding compact cutoffs
is justified and yields
\begin{equation}\label{eq:moment-source-identity}
 \E_{\nu_R}\,z\cdot\nabla\varphi_R(z)=n.
\end{equation}

We now obtain bounds independent of $R$.  Put
$m_R=\varphi_R(0)$ and $g_R=\varphi_R-m_R\ge0$.
The Hessian bound implies
$g_R(y)\le C|y|^2/2$ and $|\nabla\varphi_R(y)|\le C|y|$.
For arbitrary $y,z$, set $x=\nabla\varphi_R(z)$.  Supporting planes
at $z$ and $-z$ give the single pointwise inequality
\[
 g_R(y)\ge g_R(z)+|x\cdot y|-x\cdot z.
\]
Integrate with respect to $\nu_R$, discard the nonnegative
$\E_{\nu_R}g_R$, and use \eqref{eq:moment-source-identity}:
\[
 g_R(y)\ge\int|x\cdot y|\,d\mu_R(x)-n.
\]
On $B_1$ the density of $\mu_R$ is at least
$d_0=Z^{-1}\exp(-\max_{\overline B_1}V)>0$.
Rotational symmetry of the ball therefore bounds the integral below
by $c|y|$, where $c=d_0\int_{B_1}|x_1|\,dx>0$.  Hence
\begin{equation}\label{eq:moment-coercivity}
 c|y|-n\le g_R(y)\le\frac C2|y|^2,
 \qquad
 \left(\frac{2\pi}{C}\right)^{n/2}
 \le e^{m_R}=\int e^{-g_R(y)}\,dy
 \le e^n\int e^{-c|y|}\,dy.
\end{equation}
This bounds $m_R$ on both sides and gives one integrable exponential
majorant for all the densities $e^{-\varphi_R}$.

It remains to obtain a smooth limit and to verify that its gradient
reaches every point.  On each fixed ball, the preceding estimates
bound $\varphi_R$ and $\nabla\varphi_R$ uniformly.  Since
$Z_2\le Z_R\le Z$, equation \eqref{eq:moment-truncated} bounds
$\det D^2\varphi_R$ below by a positive constant there.
Together with $D^2\varphi_R\le CI$, this gives a uniform local bound
$\delta I\le D^2\varphi_R\le CI$: if the determinant is at least
$d>0$, every eigenvalue is at least $d/C^{n-1}$.
The right-hand sides
\[
 f_R=V(\nabla\varphi_R)+\log Z_R-\varphi_R
\]
of the logarithmic equations are uniformly Lipschitz on that ball,
because
$\nabla f_R=D^2\varphi_R\nabla V(\nabla\varphi_R)
-\nabla\varphi_R$ is uniformly bounded.

To apply an interior uniformly elliptic estimate, choose a smooth
concave scalar function $s$ equal to $\log$ on a slightly larger
interval than $[\delta,C]$, with $s'$ bounded above and bounded below
by positive constants on $\R$.  Such an extension is obtained by
smoothly continuing the decreasing derivative $1/t$ to positive
constant tails.  For symmetric matrices $Q$, the spectral function
$\mathcal F(Q)=\tr s(Q)$ is smooth, concave and uniformly elliptic:
its derivative has eigenvalues $s'(\lambda_i(Q))$, and its second
variation has coefficients $s''(\lambda_i)$ and the nonpositive
divided differences of $s'$.  It agrees with $\log\det$ on the
Hessians under consideration.  The nonhomogeneous Evans--Krylov
estimate \cite[Corollary~2.3]{Wang2012} applied to
$\mathcal F(D^2\varphi_R)=f_R$ gives uniform interior
$C^{2,\alpha}$ bounds for some $\alpha>0$.

Each $\varphi_R$ is already smooth.  Its derivative
$w_{R,k}=\partial_k\varphi_R$ satisfies
\[
 (D^2\varphi_R)^{-1}:D^2w_{R,k}
   =\nabla V(\nabla\varphi_R)\cdot\nabla w_{R,k}-w_{R,k},
\]
where $Q:N=\tr(Q^TN)$.  The principal coefficients and the
right-hand side are uniformly $C^\alpha$ on smaller balls.
Interior Schauder estimates
\cite[Theorem~2.20 and Corollary~2.21]{FernandezRealRosOton2023}
give uniform $C^{3,\alpha}$ bounds for $\varphi_R$, and iteration
gives bounds for every derivative on compact sets.
A diagonal subsequence, still indexed by $R$, converges locally smoothly to an even
$\varphi$ with positive definite Hessian, the same upper cap, and
the logarithmic equation in \eqref{eq:moment}.
The bound $\varphi_R(z)\ge c|z|-b$, with $b$ independent of $R$,
also passes to $\varphi$.

The common exponential majorant from \eqref{eq:moment-coercivity}
proves $\int e^{-\varphi}=1$ by dominated convergence.  For each
bounded continuous $h$, it also gives
\[
 \int h(\nabla\varphi)e^{-\varphi}
 =\lim_{R\to\infty}\int h(\nabla\varphi_R)e^{-\varphi_R}
 =\lim_{R\to\infty}\int h\,d\mu_R
 =\int h\,d\mu.
\]
Thus $(\nabla\varphi)_\#\nu=\mu$.
Positive definite Hessian makes the gradient locally invertible and
injective, since for $z\ne w$,
\[
 (\nabla\varphi(z)-\nabla\varphi(w))\cdot(z-w)
 =\int_0^1 D^2\varphi(w+t(z-w))[z-w,z-w]\,dt>0.
\]
Its image is dense: an open ball missed by the image would have zero
pushforward mass, contrary to the everywhere positive density of
$\mu$.  Fix $x_0\in\R^n$.  Density permits finitely many image
points $y_j=\nabla\varphi(z_j)$ whose convex hull contains a ball
$B_r(x_0)$ with $r>0$; one may take small perturbations of the
points $x_0\pm e_i$.  Their supporting planes imply
\[
 \varphi(z)-x_0\cdot z
 \ge\max_j\{(y_j-x_0)\cdot z+\varphi(z_j)-y_j\cdot z_j\}
 \ge r|z|-b_{x_0}.
\]
This function attains a minimum, at which
$\nabla\varphi=x_0$.  The gradient is therefore onto, and its local
smooth inverses combine to a global smooth inverse.
\end{proof}

We will write $x=\nabla\varphi(z)$ and
$\tau=D_z^2\varphi$, evaluating $\tau$ at the corresponding point
when using $x$ coordinates.  Thus
$\tau(x)=D_z^2\varphi((\nabla\varphi)^{-1}(x))$ is smooth, even,
positive definite and bounded above by $(4/\eps)I$.
The gradient is odd, as is its inverse.  Finally, a compactly
supported smooth function of $x$ is compactly supported after
composition with $\nabla\varphi$: its support lies in the image of
a compact set under the continuous inverse map.  Consequently the
compact integrations by parts below are valid in either coordinate
system.
